\clearpage\selectlanguage{german}\fancyhead[L]{\small ZUR THEORIE DER NEUEN TRANSSCENDENTEN. II.}\pdfbookmark[0]{Zur Theorie der transscendenten Functionen}{transscendentes-ii}\SourceLeaf{445}
\begin{center}[II.]\end{center}
\begin{center}\textbf{ZUR THEORIE DER TRANSSCENDENTEN FUNCTIONEN GEHÖRIG.}\end{center}
\par\nobreak\medskip\centerline{\rule{.14\linewidth}{.4pt}}
\TransscendentArticle{1}
Es sei $\sum e^{-\alpha(k+\omega)^2}=T$, indem für $k$ alle ganzen positiven und negativen 
Zahlen gesetzt werden, und
\[
T=A+2B\cos\omega P+2C\cos2\omega P+2D\cos3\omega P+\text{ etc.}
\]
so ist
\[
\begin{gathered}
A=\int T\,\mathrm{d}\omega,\qquad B=\int T\cos\omega P\mathbin{.}\mathrm{d}\omega,\qquad
\\[7pt]\quad{}C=\int T\cos2\omega P\mathbin{.}\mathrm{d}\omega,\\[7pt]
D=\int T\cos3\omega P\mathbin{.}\mathrm{d}\omega,\quad\text{u.\,s.\,w.}
\end{gathered}
\]
alle Integrale von $\omega=0$ bis $\omega=1$ ausgedehnt. Es ist aber klar, dass jene 
Integrale zwischen diesen Grenzen mit den Integralen
\[
\begin{gathered}\int e^{-\alpha\omega\omega}\,\mathrm{d}\omega,\quad
\int e^{-\alpha\omega\omega}\cos\omega P\mathbin{.}\mathrm{d}\omega,\quad
\\[7pt]\quad{}\int e^{-\alpha\omega\omega}\cos2\omega P\mathbin{.}\mathrm{d}\omega,\quad
\int e^{-\alpha\omega\omega}\cos3\omega P\mathbin{.}\mathrm{d}\omega,\text{ etc.}\end{gathered}
\]
übereinkommen, wenn diese von $\omega=-\infty$ bis $\omega=+\infty$ ausgedehnt werden. Da nun allgemein
\[
\begin{gathered}
e^{-\alpha\omega\omega}\cos n\omega P\\[7pt]
=\tfrac12e^{-\alpha\omega\omega+in\omega P}+\tfrac12e^{-\alpha\omega\omega-in\omega P}
\\[7pt]\quad{}=\tfrac12e^{-\frac{nnPP}{4\alpha}}
\Bigl\{e^{-\alpha\Bigl(\omega-\frac{inP}{2\alpha}\Bigr)^2}
+e^{-\alpha\Bigl(\omega+\frac{inP}{2\alpha}\Bigr)^2}\Bigr\}
\end{gathered}
\]

\SourceLeaf{446}
so folgt leicht, dass
\[
\int e^{-\alpha\omega\omega}\cos n\omega P\mathbin{.}\mathrm{d}\omega
=e^{-\frac{nn\pi\pi}{\alpha}}\surd\frac{\pi}{\alpha}
\]
folglich
\[
\begin{gathered}T=\surd\frac{\pi}{\alpha}\mathbin{.}
\Bigl\{1+2e^{-\frac{\pi\pi}{\alpha}}\cos\omega P
\\[7pt]\quad{}+2e^{-4\frac{\pi\pi}{\alpha}}\cos2\omega P
+2e^{-9\frac{\pi\pi}{\alpha}}\cos3\omega P+\text{ etc.}\Bigr\}\end{gathered}
\]
In einer andern Form so:
\[
\sum e^{-\alpha(k+\omega)^2}
=\surd\frac{\pi}{\alpha}\mathbin{.}e^{-\alpha\omega\omega}\mathbin{.}
\sum e^{-\frac{\pi\pi}{\alpha}\left(k+\frac{\alpha\omega i}{\pi}\right)^2}
\]
Allgemeiner:
\[
\begin{gathered}
\sum e^{-\alpha(k+\omega)^2}\\[7pt]\quad{}\{\cos2(k+\omega)\alpha\psi-i\sin2(k+\omega)\alpha\psi\}\\[7pt]
=\surd\frac{\pi}{\alpha}\mathbin{.}
\sum e^{-\frac{\pi\pi}{\alpha}\Bigl(k-\frac{\alpha\psi}{\pi}\Bigr)^2}
\{\cos2k\omega\pi-i\sin2k\omega\pi\}
\end{gathered}
\]
Oder, $\alpha\alpha'=\pi\pi$ und $-\alpha\psi=\omega'\pi$ gesetzt,
\[
\begin{gathered}
\sum e^{-\alpha(k+\omega)^2}\\[7pt]\quad{}\{\cos(2k+2\omega)\omega'\pi-i\sin(2k+2\omega)\omega'\pi\}\\[7pt]
=\surd\frac{\pi}{\alpha}\mathbin{.}\sum e^{-\alpha'(k+\omega')^2}
\{\cos2k\omega\pi-i\sin2k\omega\pi\}
\end{gathered}
\]
\[
\begin{gathered}\sum e^{-\alpha(k+\omega-\frac14)^2}-\sum e^{-\alpha(k+\omega+\frac14)^2}
\\[7pt]\quad{}=4\surd\frac{\pi}{\alpha}\mathbin{.}
\Bigl\{e^{-\frac{\pi\pi}{\alpha}}\sin\omega P
-3e^{-9\frac{\pi\pi}{\alpha}}\sin3\omega P+\text{ etc.}\Bigr\}\end{gathered}
\]
Man kann den Lehrsatz auch so ausdrücken:
\[
\sum t e^{-\pi t^{4}kk-2\pi ttk\surd u-\frac12\pi u}
\]
ändert den Werth nicht, wenn $t$ in $\dfrac1t$ und $u$ in $-u$ verwandelt wird.
\TransscendentArticle{2}
Man setze
\[
\begin{gathered}P=1+\frac{x^n\mathbin{.}1-x^{2n+1}}{1+x^n\mathbin{.}1+x^{n+1}}
\\[7pt]\quad{}+\frac{x^{2n}\mathbin{.}1-x^{2n+2}\mathbin{.}1-x^{n+1}}{1+x^n\mathbin{.}1+x^{n+1}\mathbin{.}1+x^{n+2}}
\\[7pt]\quad{}+\frac{x^{3n}\mathbin{.}1-x^{2n+3}\mathbin{.}1-x^{n+1}\mathbin{.}1-x^{n+2}}{1+x^n\mathbin{.}1+x^{n+1}\mathbin{.}1+x^{n+2}\mathbin{.}1+x^{n+3}}
+\text{ etc.}\end{gathered}
\]
\[
\begin{gathered}Q=\frac{x^n}{1+x^n}
+\frac{x^{2n}\mathbin{.}1-x^{n+1}}{1+x^n\mathbin{.}1+x^{n+1}}
\\[7pt]\quad{}+\frac{x^{3n}\mathbin{.}1-x^{n+1}\mathbin{.}1-x^{n+2}}{1+x^n\mathbin{.}1+x^{n+1}\mathbin{.}1+x^{n+2}}
+\text{ etc.}\end{gathered}
\]
\[
R=P-Q
\]

\SourceLeaf{447}
Man suche zuerst diese Differenz, indem man das erste, zweite, dritte Glied 
u.\,s.\,w. der Reihe $Q$ von dem ersten, zweiten, dritten Gliede der Reihe $P$ abzieht, so kommt
\[
\begin{gathered}R=\frac1{1+x^n}
\\[7pt]\quad{}+\frac{x^n\mathbin{.}1-x^n}{1+x^n\mathbin{.}1+x^{n+1}}
\\[7pt]\quad{}+\frac{x^{2n}\mathbin{.}1-x^n\mathbin{.}1-x^{n+1}}{1+x^n\mathbin{.}1+x^{n+1}\mathbin{.}1+x^{n+2}}
\\[7pt]\quad{}+\frac{x^{3n}\mathbin{.}1-x^n\mathbin{.}1-x^{n+1}\mathbin{.}1-x^{n+2}}{1+x^n\mathbin{.}1+x^{n+1}\mathbin{.}1+x^{n+2}\mathbin{.}1+x^{n+3}}
+\text{etc.}\end{gathered}
\]
Wir bezeichnen diese Reihe durch $\varphi(x,n)$.

Man suche ferner jene Differenz $R$, indem man das erste, zweite, dritte 
Glied u.\,s.\,w. der Reihe $Q$ von dem zweiten, dritten, vierten u.\,s.\,w. der Reihe 
$P$ abzieht, so wird
\[
\begin{gathered}R=1-\frac{x^{2n+1}}{1+x^{n+1}}
\\[7pt]\quad{}-\frac{x^{3n+2}\mathbin{.}1-x^{n+1}}{1+x^{n+1}\mathbin{.}1+x^{n+2}}
\\[7pt]\quad{}-\frac{x^{4n+3}\mathbin{.}1-x^{n+1}\mathbin{.}1-x^{n+2}}{1+x^{n+1}\mathbin{.}1+x^{n+2}\mathbin{.}1+x^{n+3}}
-\text{ u.\,s.\,w.}\end{gathered}
\]
oder offenbar
\[
R=1-x^{2n+1}\mathbin{.}\varphi(x,n+1)
\]
folglich
\[
\varphi(x,n)=1-x^{2n+1}\mathbin{.}\varphi(x,n+1)
\]
Dieser Schluss ist allgemein, so lange $n>1$, man hat demnach unter dieser Einschränkung
\[
\varphi(x,n)=1-x^{2n+1}+x^{4n+4}-x^{6n+9}+x^{8n+16}-\text{ etc.}
\]
hingegen ist für den Fall $n=0$ das letzte Glied von $Q$ nicht als verschwindend zu betrachten. Setzt man es $=T$, so wird der erste Werth von $R$ um 
$T$ kleiner sein als der zweite, also
\[
T=1-x\varphi(x,1)-\varphi(x,0)
\]
oder
\[
\begin{gathered}\varphi(x,0)=1-x\varphi(x,1)-T\\[7pt]\quad{}=1-x+x^4-x^9+x^{16}-\ .\ .\ -T\end{gathered}
\]
also da
\[
\varphi(x,0)=\tfrac12
\]
und
\[
2T=\frac{1-x}{1+x}\mathbin{.}\frac{1-xx}{1+xx}\mathbin{.}\frac{1-x^3}{1+x^3}\mathbin{.}\frac{1-x^4}{1+x^4}\quad\text{etc.}
\]
so wird
\SourceLeaf{448}
\[
\begin{gathered}\frac{1-x}{1+x}\,.\,\frac{1-xx}{1+xx}\,.\,\frac{1-x^3}{1+x^3}\,.\,\frac{1-x^4}{1+x^4}\ \text{etc.}
\\[7pt]\quad{}=1-2x+2x^4-2x^9+2x^{16}-\text{etc.}\end{gathered}
\]
der erste Theil ist hier
\[
\begin{gathered}
=1-x\,.\,1-x^3\,.\,1-x^5\,.\,1-x^7\ \text{etc.}\quad
\\[7pt]\quad{}\frac{1-xx}{1+x}\,.\,\frac{1-x^4}{1+xx}\,.\,\frac{1-x^6}{1+x^3}\ \text{etc.}\\[7pt]
=(1-x)^2(1-xx)(1-x^3)^2(1-x^4)(1-x^5)^2(1-x^6)\ \text{etc.}
\end{gathered}
\]
Bezeichnen wir noch
\[
1-2x+2x^4-2x^9+2x^{16}\ \text{etc.}\quad\text{mit }Fx
\]
so wird offenbar
\[
Fx\,.\,F(-x)=(Fxx)^2
\]
\TransscendentArticle{3}
Zieht man auf ähnliche Weise von der Reihe
\[
\begin{gathered}\frac{1-x^{2n+2}}{1-x^{n+1}}
\\[7pt]\quad{}+\frac{x^n\,.\,1-x^{2n+4}\,.\,1-x^{n+2}}{1-x^{n+1}\,.\,1-x^{n+3}}
\\[7pt]\quad{}+\frac{x^{2n}\,.\,1-x^{2n+6}\,.\,1-x^{n+2}\,.\,1-x^{n+4}}{1-x^{n+1}\,.\,1-x^{n+3}\,.\,1-x^{n+5}}
+\text{etc.}\end{gathered}
\]
die Reihe
\[
\begin{gathered}\frac{x^n\,.\,1-x^{n+2}}{1-x^{n+1}}
\\[7pt]\quad{}+\frac{x^{2n}\,.\,1-x^{n+2}\,.\,1-x^{n+4}}{1-x^{n+1}\,.\,1-x^{n+3}}
\\[7pt]\quad{}+\frac{x^{3n}\,.\,1-x^{n+2}\,.\,1-x^{n+4}\,.\,1-x^{n+6}}{1-x^{n+1}\,.\,1-x^{n+3}\,.\,1-x^{n+5}}
+\text{etc.}\end{gathered}
\]
ab, so erhält man einmal
\[
\begin{gathered}\frac{1-x^n}{1-x^{n+1}}
\\[7pt]\quad{}+\frac{x^n\,.\,1-x^n\,.\,1-x^{n+2}}{1-x^{n+1}\,.\,1-x^{n+3}}
\\[7pt]\quad{}+\frac{x^{2n}\,.\,1-x^n\,.\,1-x^{n+2}\,.\,1-x^{n+4}}{1-x^{n+1}\,.\,1-x^{n+3}\,.\,1-x^{n+5}}
+\text{etc.}=\psi(x,n)\end{gathered}
\]
und zweitens
\[
\begin{gathered}
1+x^{n+1}+\frac{x^{2n+3}\,.\,1-x^{n+2}}{1-x^{n+3}}
\\[7pt]\quad{}+\frac{x^{2n+3}\,.\,x^{n+2}\,.\,1-x^{n+2}\,.\,1-x^{n+4}}{1-x^{n+3}\,.\,1-x^{n+5}}+\text{etc.}\\[7pt]
=1+x^{n+1}+x^{2n+3}\psi(x,n+2)
\end{gathered}
\]
Man hat also, den Fall \(n=0\) ausgenommen,
\[
\psi(x,n)=1+x^{n+1}+x^{2n+3}\psi(x,n+2)
\]
folglich
\[
\psi(x,n)=1+x^{n+1}+x^{2n+3}+x^{3n+6}+x^{4n+10}+\text{etc.}
\]
dagegen hat man für den Fall \(n=0\)
\[
\begin{gathered}\psi(x,0)=1+x+x^3\psi(x,2)\\[7pt]\quad{}-\frac{1-x^2\,.\,1-x^4\,.\,1-x^6}{1-x\,.\,1-x^3\,.\,1-x^5}\ \text{etc.}\end{gathered}
\]

\SourceLeaf{449}
folglich
\[
\frac{1-x^2\,.\,1-x^4\,.\,1-x^6\ .\ .\ .}{1-x\,.\,1-x^3\,.\,1-x^5\ .\ .\ .}
=1+x+x^3+x^6+x^{10}+\text{etc.}
\]
\TransscendentArticle{4}
Wir bezeichnen
\[
1-x\,.\,1-xx\,.\,1-x^3\ \text{etc.}\qquad\text{mit }[x]
\]
so ist:
\[
\begin{gathered}
1+x+x^3+x^6+\text{etc.}
\\[7pt]\quad{}=\frac{1-xx\,.\,1-x^4\,.\,1-x^6\ .\ .\ .}{1-x\,.\,1-x^3\,.\,1-x^5\ .\ .\ .}\\[7pt]\quad{}=\frac{[xx]^2}{[x]}\\[7pt]
=1+x\,.\,1+xx\,.\,1+x^3\,.\,1+x^4\,.\,1+x^5\ .\ .\ .\\[7pt]\quad{}1-x^2\,.\,1-x^4\,.\,1-x^{16}\ .\ .\ .
\end{gathered}
\]
\[
\begin{gathered}1-2x+2x^4-2x^9+\text{etc.}
\\[7pt]\quad{}=\frac{1-x\,.\,1-xx\,.\,1-x^3\ .\ .\ .}{1+x\,.\,1+xx\,.\,1+x^3\ .\ .\ .}\\[7pt]\quad{}=\frac{[x]^2}{[xx]}\end{gathered}
\]
\[
\begin{gathered}
1+2x+2x^4+2x^9+\text{etc.}
\\[7pt]\quad{}=\frac{1+x\,.\,1-xx\,.\,1+x^3\ .\ .\ .}{1-x\,.\,1+xx\,.\,1-x^3\ .\ .\ .}\\[7pt]\quad{}=\frac{[xx]^5}{[x]^2[x^4]^2}\\[7pt]
=(1+x)^2(1-xx)(1+x^3)^2(1-x^4)(1+x^5)^2\ .\ .\ .
\end{gathered}
\]
\[
[-x]=\frac{[xx]^3}{[x][x^4]}
\]
\[
\begin{gathered}
(1+xy)(1+x^3y)(1+x^5y)\ .\ .\ .
\\[7pt]\quad{}\Bigl(1+\frac{x}{y}\Bigr)\Bigl(1+\frac{x^3}{y}\Bigr)\Bigl(1+\frac{x^5}{y}\Bigr)\ .\ .\ .\\[7pt]
\qquad=\frac{1}{[xx]}\Bigl\{1+x\Bigl(y+\frac{1}{y}\Bigr)+x^4\Bigl(yy+\frac{1}{yy}\Bigr)+x^9\Bigl(y^3+\frac{1}{y^3}\Bigr)+\,.\Bigr\}
\end{gathered}
\]
\[
\begin{gathered}\text{etc.}+x^{(\omega-1)^2}+x^{\omega\omega}+x^{(\omega+1)^2}+\text{etc.}
\\[7pt]\quad{}=[xx]\frac{\ .\ .\ .1+x^{2\omega+3}\,.\,1+x^{2\omega+1}\,.\,1+x^{2\omega-1}\,.\,1+x^{2\omega-3}\ .\ .\ .}{x^{-\omega\omega}\,.\,x^{2\omega-1}\,.\,x^{2\omega-3}\ .\ .\ .}\end{gathered}
\]
\[
[x]^3=1-3x+5x^3-7x^6+9x^{10}-\text{etc.}
\]
folgt leicht aus
\[
\begin{gathered}
\Bigl\{\Bigl(y-\frac{1}{y}\Bigr)x-\Bigl(y^3-\frac{1}{y^3}\Bigr)x^9+\Bigl(y^5-\frac{1}{y^5}\Bigr)x^{25}-\,.\,.\Bigr\}
\\[7pt]\quad{}\Bigl\{\Bigl(y+\frac{1}{y}\Bigr)x+\Bigl(y^3+\frac{1}{y^3}\Bigr)x^9+\Bigl(y^5+\frac{1}{y^5}\Bigr)x^{25}+\,.\,.\Bigr\}\\[7pt]
\qquad=\Bigl\{1-2x^8+2x^{32}-\,.\,.\Bigr\}
\\[7pt]\quad{}\Bigl\{\Bigl(yy-\frac{1}{yy}\Bigr)xx-\Bigl(y^6-\frac{1}{y^6}\Bigr)x^{18}+\text{etc.}\Bigr\}
\end{gathered}
\]
wenn man \(y=1+\omega\) setzt und daraus die Bedingungsgleichungen bildet.
\SourceLeaf{450}
\TransscendentArticle{5}
Zu den Hauptsätzen in dieser Theorie gehört folgendes Theorem. 
Bezeichnet man
\[
\begin{gathered}\{e^{-\frac{P\alpha}{24}}-e^{-\frac{25P\alpha}{24}}
-e^{-\frac{49P\alpha}{24}}\\[7pt]\quad{}+e^{-\frac{121P\alpha}{24}}
+e^{-\frac{169P\alpha}{24}}+\mathinner{..}\}\mathbin{.}{}^{4}\!\surd\alpha
\quad\text{durch }\varphi\alpha\end{gathered}
\]
so ist
\[
\varphi\alpha=\varphi\frac{1}{\alpha}
\]
Diese Function ist ein Maximum für $\alpha=1$, wo ihr Werth
\[
=0{,}7682255=\frac{1}{{}^{4}\!\surd2\mathbin{.}\surd1{,}19814\mathinner{..}}
={}^{4}\!\surd\tfrac12\mathbin{.}\surd\frac{\varpi}{\pi}
\]
ferner ist
\[
(\varphi 2\lambda)^{24}=4(\varphi\lambda)^8(\varphi 4\lambda)^8
\{(\varphi\lambda)^8+(\varphi 4\lambda)^8\}
\]
\[
\varphi\alpha=e^{-\frac{\alpha}{24}P}{}^{4}\!\surd\alpha\mathbin{.}
(1-e^{-aP})(1-e^{-2\alpha P})(1-e^{-3\alpha P})\ .\ .\ .
\]
oder
\[
[x]=\frac{\varphi\frac{-\log x}{P}\mathbin{.}{}^{24}\!\surd\frac1x}
{{}^{4}\!\surd\frac{-\log x}{P}}
\]

Es sei
\[
\begin{gathered}(1+xy)(1+x^3y)(1+x^5y)\ .\ .\ .
\\[7pt]\quad{}(1+\tfrac{x}{y})(1+\tfrac{x^3}{y})(1+\tfrac{x^5}{y})\ .\ .\ .=F(x,y)\end{gathered}
\]
so ist
\[
\begin{gathered}F\bigl[e^{-\frac12\alpha P},\quad e^{P\omega\surd\alpha}\bigr]
\mathbin{.}e^{-\frac{P}{24\alpha}}
\\[7pt]\quad{}=F\bigl[e^{-\frac12\frac{P}{\alpha}},\quad
e^{iP\omega\surd\frac1\alpha}\bigr]
\mathbin{.}e^{-\frac{P\alpha}{24}}\mathbin{.}e^{\frac12P\omega\omega},\end{gathered}
\]

\TransscendentArticle{6}
Die Siebentheilung führt auf folgende Gleichung
\[
\begin{gathered}
b=\Bigl[\frac{1-2x+2x^4-\mathinner{..}}{1+2x+2x^4+\mathinner{..}}\Bigr]^2\\[7pt]
A=\Bigl[\frac{1+2x^7+2x^{28}+\mathinner{..}}{1+2x+2x^4+\mathinner{..}}\Bigr]^2\\[7pt]
B=\Bigl[\frac{1-2x^7+2x^{28}+\mathinner{..}}{1+2x+2x^4+\mathinner{..}}\Bigr]^2
\end{gathered}
\]

\SourceLeaf{451}
\[
\begin{gathered}
A^8-\tfrac47A^6+\frac{16-32bb}{49}A^5-\frac{30}{343}A^4
\\[7pt]\quad{}+\frac{32-64bb}{2401}A^3
\\[7pt]\quad{}-\frac{20+768bb-768b^4}{16807}A^2\\[7pt]
+\frac{48-2144bb+6144b^4-4096\mathbin{.}b^6}{823543}A
-\frac1{823543}=0
\end{gathered}
\]
\[
\begin{gathered}
B^8-\tfrac47bbB^6+\frac{16b^3-32b}{49}B^5-\frac{30b^4}{343}B^4
\\[7pt]\quad{}+\frac{32b^5-64b^3}{2401}B^3
\\[7pt]\quad{}-\frac{20b^6+768b^4-768b^2}{16807}B^2\\[7pt]
+\frac{48b^7-2144b^5+6144b^3-4096b}{823543}B
-\frac1{823543}=0
\end{gathered}
\]
Wenn man in der ersten Gleichung statt $bb$, $1-bb$ und statt $A$, $-A$ 
setzt, so wird sie nicht geändert.

\TransscendentArticle{7}
Für die Dreitheilung hat man
\[
\begin{array}{ccc}
a&b&\dfrac ba=t\\[5pt]
A&B&\dfrac BA=T
\end{array}
\]
\[
\begin{gathered}T^4+(12t-16t^3)T^3+6ttT^3\\[7pt]\quad{}-(16t-12t^3)T+t^4=0\end{gathered}
\]
oder
\[
(T-t)^4=16(t-t^3)(T-T^3)
\]
\[
\frac Aa=\frac{3t-t^3+T(1-3tt)}{3[T(1-tt)+t(1-TT)]}
\]
Setzt man $T=\operatorname{tang}N$, $t=\operatorname{tang}n$, so ist
\[
\begin{gathered}
3\sin(2N-2n)^2=4\sin(N+3n)\mathbin{.}\sin(3N+n)\\[7pt]
\sin(N-n)^4=\sin4n\mathbin{.}\sin4N
\end{gathered}
\]
\[
\begin{gathered}\frac{A\cos n}{a\cos N}=\frac{B\sin n}{b\sin N}
\\[7pt]\quad{}=\frac{2\sin(N+3n)}{3\sin(2N-2n)}
=\frac{\sin(2N-2n)}{2\sin(3N+n)}\end{gathered}
\]

\TransscendentArticle{8}
Zum Beweise der schönen Lehrsätze der Reciprocität wird folgendes dienen: 
I.\quad Das Product aus allen
\[
1-\frac{\eta}{ak+N}
\]
ist das Product aus allen
